A continuación se responden las diez preguntas de análisis final que pide la guía, con base en los
resultados experimentales obtenidos a lo largo de todas las actividades.

\begin{enumerate}

\item \textbf{¿Cuál es la diferencia fundamental entre UCS y A*?}

UCS selecciona para expansión el estado con menor $f(n) = g(n)$, es decir, únicamente el costo
acumulado real desde el inicio. A* selecciona el estado con menor $f(n) = g(n) + h(n)$,
incorporando además una estimación heurística del costo restante hasta la meta. Esa información
adicional es lo que permite a A* orientar la búsqueda hacia la meta y expandir menos estados que
UCS, sin perder la garantía de optimalidad siempre que $h(n)$ sea admisible.

\item \textbf{¿Qué función cumple $g(n)$?}

$g(n)$ es el costo acumulado real (conocido, no estimado) del camino recorrido desde el estado
inicial hasta el estado $n$. Es la misma cantidad que UCS usa como única prioridad.

\item \textbf{¿Qué función cumple $h(n)$?}

$h(n)$ es una estimación del costo que falta recorrer desde $n$ hasta un estado objetivo. No es un
valor conocido con exactitud (salvo en el caso trivial $h(n)=0$): es una heurística, calculada con
información parcial del problema (por ejemplo, distancia Manhattan, distancia Euclidiana, o el
diámetro de posiciones pendientes por visitar en los problemas de esquinas y comida).

\item \textbf{¿Qué ocurre cuando $h(n) = 0$?}

A* se reduce exactamente a UCS, porque $f(n) = g(n) + 0 = g(n)$: la prioridad usada para elegir el
siguiente estado a expandir es idéntica en ambos algoritmos. Esto se confirmó empíricamente en dos
problemas distintos: en \texttt{mediumClassic}, UCS y A*+h(n)=0 expanden exactamente 69 nodos cada
uno (Actividad 4); en \texttt{tinyCorners}, UCS y A*+h(n)=0 expanden exactamente 295 nodos cada uno
(Actividad 9). En ambos casos el costo óptimo encontrado también es idéntico.

\item \textbf{¿Cuál presentó mejores resultados: Manhattan o Euclidiana?}

Manhattan, con menos nodos expandidos (15 contra 16 en \texttt{mediumClassic}, Actividad 6), aunque
ambas heurísticas encuentran el mismo costo óptimo (12). La razón es que Pac-Man solo se mueve en
las cuatro direcciones cardinales (nunca en diagonal), y matemáticamente $h_E(n) \le h_M(n)$
siempre (la distancia en línea recta nunca es mayor que la distancia recorrida por una cuadrícula).
Por lo tanto Manhattan es una cota más ajustada --más informada-- sin dejar de ser admisible, y una
heurística admisible más informada expande menos nodos (ver pregunta 6).

\item \textbf{¿Una heurística más grande siempre es mejor? Justifique.}

Solo si sigue siendo admisible, es decir, si sigue cumpliendo $h(n) \le h^*(n)$. Entre dos
heurísticas admisibles, la que devuelve valores mayores (más informada) expande menos o igual
número de nodos que la menos informada, sin perder optimalidad --esto se demostró repetidamente en
este laboratorio: en el problema de esquinas, la heurística propuesta (diámetro) siempre devuelve
un valor mayor o igual que la básica, y expande menos nodos (119 contra 147, Actividad 8); en el
problema de comida, la Heurística 2 (diámetro) expande menos que la Heurística 1 (383 contra 702 en
\texttt{testClassic}, Actividad 11). Si la heurística dejara de ser admisible (sobreestimara en
algún estado), ``más grande'' dejaría de ser sinónimo de ``mejor'' --véase la pregunta 7--.

\item \textbf{¿Por qué una heurística que sobreestima puede ser problemática?}

Porque rompe la condición $h(n) \le h^*(n)$ que garantiza la optimalidad de A*. Si en algún estado
$h(n)$ es mayor que el costo real restante, A* puede descartar --o posponer indefinidamente-- un
camino que en realidad era parte de la solución óptima, porque su $f(n) = g(n) + h(n)$ aparenta ser
peor de lo que realmente es. El algoritmo puede terminar devolviendo una solución de costo mayor al
óptimo sin que nada en su ejecución indique que ocurrió un error: el resultado se ve igual de
``válido'' que uno correcto, pero no lo es.

\item \textbf{¿Por qué el estado de \texttt{CornersProblem} necesita almacenar más información que
únicamente $(x, y)$?}

Porque la prueba de objetivo de este problema no depende solo de dónde está Pac-Man, sino de qué ha
hecho para llegar ahí --específicamente, cuáles de las 4 esquinas ya visitó--. Si el estado fuera
solo la posición, dos situaciones con la misma posición pero distinto progreso (unas esquinas
visitadas y otras no) se confundirían en un único estado, y un algoritmo de búsqueda podría
descartar como ``ya explorado'' un estado que en realidad todavía necesitaba explorarse con otra
combinación de esquinas pendientes, perdiendo la garantía de encontrar la solución óptima --o
incluso de encontrar alguna solución-- (desarrollado en detalle en la Actividad 7).

\item \textbf{¿Por qué \texttt{FoodSearchProblem} presenta un espacio de estados considerablemente
mayor?}

Porque el estado incluye qué alimentos siguen presentes, y cada uno de los $F$ alimentos puede
encontrarse en dos condiciones (\{presente, consumido\}), por lo que el número de configuraciones
posibles de comida crece como $2^F$, multiplicado además por cada posición posible de Pac-Man. Esto
se confirmó empíricamente: con apenas 8 alimentos (\texttt{testClassic}), UCS/A*+h=0 ya expanden
2598 nodos (Actividad 10); con 55 alimentos (\texttt{smallClassic}) la búsqueda no terminó en 45-60
segundos, ilustrando de forma directa la explosión combinatoria de $2^F$.

\item \textbf{¿Qué relación encontró entre la calidad de la heurística y el número de nodos
expandidos?}

Una relación inversa y consistente en los tres problemas trabajados: entre más informada (más
cercana al costo real, sin dejar de ser una cota inferior) es la heurística, menos nodos necesita
expandir A* para encontrar la misma solución óptima. Se observó la misma progresión en dos
problemas distintos: en esquinas, $h=0 \to$ básica $\to$ propuesta redujo la exploración de 295 a
147 a 119 nodos (Actividad 9, factor de reducción $R=2.48$ con la propuesta); en comida, $h=0 \to$
Heurística 1 $\to$ Heurística 2 redujo de 2598 a 702 a 383 nodos sobre \texttt{testClassic}
(Actividad 11). En ambos casos el costo óptimo permaneció idéntico en todas las estrategias,
confirmando que la reducción de exploración no tiene ningún costo en la calidad de la solución
--siempre que la heurística conserve la admisibilidad--.

\end{enumerate}
